




% Table generated by Excel2LaTeX from sheet 'Sheet1'
\begin{table*}[t]
	\centering
	%	\vspace{-8pt}
	\caption{A coarse summary of IL methods demonstrating 1) different expert data modalities they can handle (learning from {\textit{demonstrations}} or {\textit{observations}}), 2) disparate task settings they consider (learning from {\textit{online}} environment interactions or pre-collected {\textit{offline}} static dataset), 3) the specific {\textit{cross-domain}} setting they assume (the transition dynamics between the learning environment and that of the expert behaviors are different), and 4) the unique {\textit{one-shot}} merit they desire (the learned policy is capable of one-shot transfer to new imitation tasks). 
    We highlight that our contextual imitation learning (CEIL) method can naturally be applied to all the above IL settings. % without introducing any additional models or losses. 
	}
	\vspace{5pt}
		\resizebox{\linewidth}{!}{
		\begin{tabular}{lccccccccc}
			\toprule
			& \multicolumn{2}{c}{Expert data} & \multicolumn{2}{c}{Task setting} & \multirow{2}{*}{\makecell[c]{Cross-domain}} & \multirow{2}{*}{One-shot} \\
			\cmidrule(lr){2-3} \cmidrule(lr){4-5} 
			& \multicolumn{1}{c}{~~LfD~~~} & \multicolumn{1}{c}{~~LfO~~~} & \multicolumn{1}{c}{Online} & \multicolumn{1}{c}{Offline} \\% & \multicolumn{1}{c}{ } &   \\
			\midrule
            S-on-LfD \citep{ brown2020better, dadashi2020primal, fu2017learning,ho2016generative,ke2021imitation, liu2020energy, ni2021f, reddy2019sqil, ziebart2008maximum}  
			& \mineyes & \mineno & \mineyes & \mineno & \mineno  & \mineno \\
            S-on-LfO \citep{boborzi2022imitation, liu2018imitation, torabi2018behavioral, torabi2018generative, zhu2020off}  
			& \mineyes & \mineyes & \mineyes & \mineno & \mineno  & \mineno \\
            S-off-LfD \citep{florence2022implicit, jarboui2021offline, jarrett2020strictly, kim2022demodice, luo2023optimal, xu2022discriminator, yue2023clare, zhang2023discriminator} 
            & \mineyes & \mineno & \mineno & \mineyes & \mineno  & \mineno \\
            S-off-LfO \citep{zolna2020offline}
            & \mineyes & \mineyes & \mineno & \mineyes & \mineno & \mineno \\
            C-on-LfD \citep{fickinger2021cross, wang2022latent, zolna2021task} 
            & \mineyes & \mineno & \mineyes & \mineno & \mineyes & \mineno  \\
            C-on-LfO \citep{franzmeyer2022learn, gangwani2020state, gangwani2022imitation, liu2019state, qiu2023out, raychaudhuri2021cross}
			& \mineyes & \mineyes & \mineyes & \mineno & \mineyes & \mineno  \\
            C-off-LfD \citep{jiang2020offline}
            & \mineyes & \mineno & \mineno & \mineyes & \mineyes & \mineno  \\
            C-off-LfO \citep{ma2022smodice, viano2022robust}
			& \mineyes & \mineyes & \mineno & \mineyes & \mineyes & \mineno \\
            \midrule
            S-on/off-LfO \citep{garg2021iq}
			& \mineyes & \mineyes & \mineyes & \mineyes & \mineno & \mineno \\
            Online one-shot \citep{dance2021conditioned, duan2017one, kim2020domain}
            & \mineyes & \mineno & \mineyes & \mineno & \mineno & \mineyes \\
            Offline one-shot \citep{furuta2021generalized, xu2022prompting}
            & \mineyes & \mineno & \mineno & \mineyes & \mineno & \mineyes \\

			% ValueDICE~\citep{kostrikov2019imitation}
			% & \mineyes & \mineno & \mineyes & \mineno & \mineno  \\
   
			
			\toprule
			\textbf{CEIL (ours)}  
			& \mineyes & \mineyes & \mineyes & \mineyes & \mineyes & \mineyes  \\
			\bottomrule
		\end{tabular}%
			}
	\label{tab:ceil-on-various-IL-settings}%
	%	\vspace{-9pt}
\end{table*}%% \mineno    & {\scriptsize \CheckmarkBold  }








\section{Related Work}

Recent advances in decision-making have led to rapid progress in IL settings (Table~\ref{tab:ceil-on-various-IL-settings}), from typical learning from demonstrations (LfD) to learning from observations (LfO)~\citep{boborzi2022imitation, brown2020better, judah2014imitation, liu2018imitation, ross2011reduction, torabi2018generative}, from online IL to offline IL~\citep{chang2021mitigating, demoss2023ditto, jarrett2020strictly,  liu2021curriculum,zhang2023discriminator}, and from single-domain IL to cross-domain IL~\citep{jiang2020offline, liu2019state, ma2022smodice, viano2022robust}. 
Targeting a specific IL setting, individual works have shown their impressive ability to solve the exact IL setting. However, it is hard to retrain their performance in new unprepared IL settings. In light of this, it is tempting to consider how we can design a general and broadly applicable IL method. Indeed, a number of prior works have studied part of the above IL settings, such as offline LfO~\citep{zolna2020offline}, cross-domain LfO~\citep{liu2019state, raychaudhuri2021cross}, and cross-domain offline IL~\citep{ma2022smodice}. 
While such works demonstrate the feasibility of tackling multiple IL settings, they still rely on standard online/offline RL algorithmic advances to improve performance~\citep{gangwani2020state, jarboui2021offline, lee2021generalizable, li2017infogail, liu2022dara, liu2023beyond, luo2023optimal, yue2023clare, zhuang2023behavior}. 
% The objective of our method 
Our objective diverges from these works, as we strive to minimize the reliance on the RL pipeline by replacing it with a simple supervision objective, thus avoiding the dependence on the choice of RL algorithms.


Our approach to IL is most closely related to prior hindsight information-matching methods~\citep{ajay2022conditional, brandfonbrener2022does, furuta2021generalized, liu2022learn}, both learning a contextual policy and using a contextual variable to guide policy improvement. 
However, these prior methods typically require additional mechanisms to work well, such as extrinsic rewards in online RL~\citep{arulkumaran2022all, kumar2019reward, srivastava2019training} or a handcrafted target return in offline RL~\citep{chen2021decision, emmons2021rvs}. Our method does not require explicit handling of these components. 
By explicitly learning an embedding space for both expert and suboptimal behaviors, we can bias the contextual policy with an inferred optimal embedding (contextual variable), thus avoiding the need for explicit reward biasing in prior works. 
Our method also differs from most prior offline transformer-based RL/IL algorithms that explicitly model a long sequence of transitions~\citep{carroll2022unimask, chen2021decision, janner2021offline, kang2023beyond, lai2023chipformer, xu2022prompting}. We find that simple fully-connected networks can also elicit useful embeddings and guide expert behaviors when conditioned on a well-calibrated embedding. 
In the context of the recently proposed prompt-tuning paradigm in large language tasks or multi-modal tasks~\citep{gao2021clip, lester2021power, zhou2022learning}, our method can be interpreted as a combination of IL and prompting-tuning, with the main motivation that we tune the prompt (the optimal contextual variable) with an expert matching objective in IL settings. 

% contrastive loss while most prompt-tuning methods use the supervision loss to tune it. 









